\section{Exact singular-value geometry of a fiber sum}
\label{sec:fiber-svd}

The terminal arithmetic will be inserted only after the finite-dimensional
operator has been solved exactly.

\subsection{Weighted atomic spaces}

Let \(\cU\) and \(\cY\) be finite sets, let
\(q:\cU\to\cY\), and assign positive atomic weights \(a_u>0\) and
\(b_y>0\).  Write
\begin{equation}
 \|z\|_a^2=\sum_{u\in\cU}a_u|z_u|^2,
 \qquad
 \|w\|_b^2=\sum_{y\in\cY}b_y|w_y|^2.
 \label{eq:tpc55-weighted-norms}
\end{equation}
For coefficients \(c_u\in\C\), define
\begin{equation}
 (Tz)_y=\sum_{u\in F_y}c_uz_u,
 \qquad F_y=q^{-1}(y).
 \label{eq:tpc55-general-fiber-operator}
\end{equation}
Put
\begin{equation}
 S_y:=\sum_{u\in F_y}\frac{|c_u|^2}{a_u},
 \qquad
 \rho_y^2:=b_yS_y,
 \qquad
 \cY_+:=\{y:S_y>0\}.
 \label{eq:tpc55-fiber-row-norm}
\end{equation}

\begin{theorem}[Complete weighted fiber SVD]
\label{thm:tpc55-complete-svd}
The multiset of nonzero singular values of \(T\) is exactly
\begin{equation}
 \{\!\{\rho_y:y\in\cY_+\}\!\},
\label{eq:tpc55-all-singular-values}
\end{equation}
with one occurrence contributed by each active output fiber.  Thus distinct
fibers with the same value of \(\rho_y\) contribute the corresponding
singular-value multiplicity.  In particular,
\begin{equation}
 \rank T=|\cY_+|,
 \qquad
 \dim\ker T=|\cU|-|\cY_+|.
 \label{eq:tpc55-rank-nullity}
\end{equation}
For \(y\in\cY_+\), the corresponding unit right and left singular
vectors are
\begin{equation}
 (v_y)_u=
 \begin{cases}
 \dfrac{\overline{c_u}}{a_u\sqrt{S_y}},&u\in F_y,\\[6pt]
 0,&u\notin F_y,
 \end{cases}
 \qquad
 f_y=b_y^{-1/2}e_y,
 \label{eq:tpc55-singular-vectors}
\end{equation}
and \(Tv_y=\rho_yf_y\).

If \(P_y\) is the \(a\)-orthogonal projection onto \(\C v_y\), then
\begin{equation}
 (P_yz)_u
 =\frac{\overline{c_u}}{a_uS_y}
   \sum_{v\in F_y}c_vz_v
 \quad(u\in F_y),
 \label{eq:tpc55-visible-projection}
\end{equation}
and, with \(P=\sum_{y\in\cY_+}P_y\),
\begin{equation}
 \cH_a(\cU)
 =\ker T\mathbin{\mathop{\oplus}\limits^{\perp_a}}
  \operatorname{Ran}P,
 \qquad
 T^*T=\sum_{y\in\cY_+}\rho_y^2P_y.
 \label{eq:tpc55-svd-decomposition}
\end{equation}
\end{theorem}

\begin{proof}
Conjugate \(T\) by the isometries
\(x_u=\sqrt{a_u}z_u\) and \(W_y=\sqrt{b_y}w_y\).  The resulting matrix
has row
\[
 d_{y,u}=\sqrt{\frac{b_y}{a_u}}c_u\one_{q(u)=y}.
\]
Distinct rows have disjoint supports, hence are orthogonal, and the squared
norm of row \(y\) is \(b_yS_y=\rho_y^2\).  Normalizing the nonzero rows
gives \eqref{eq:tpc55-singular-vectors}.  The projection formula follows
by taking the \(a\)-inner product with \(v_y\).  The remaining identities
are the spectral decomposition and rank--nullity theorem.
\end{proof}

Thus a collision fiber has no complicated internal spectrum: it has one
visible direction and a completely flat zero spectrum in every transverse
direction.

\subsection{Kernel basis and coordinate transversals}

\begin{proposition}[Explicit collision kernel]
\label{prop:tpc55-kernel-basis}
Fix \(y\in\cY_+\) and choose \(u_0\in F_y\) with \(c_{u_0}\ne0\).  Then
\begin{equation}
 k_{y,u}:=c_ue_{u_0}-c_{u_0}e_u,
 \qquad u\in F_y\setminus\{u_0\},
 \label{eq:tpc55-kernel-vectors}
\end{equation}
is a basis of the kernel of the fiber functional
\(z\mapsto\sum_{u\in F_y}c_uz_u\).  Hence a coordinate subspace
supported on \(S\subseteq\cU\) is a domain of injectivity for \(T\) if
and only if
\begin{equation}
 c_u\ne0\quad(u\in S),
 \qquad
 q|_S\ \text{is injective}.
 \label{eq:tpc55-coordinate-injectivity}
\end{equation}
When these conditions hold, the singular values of the restriction are
exactly
\begin{equation}
 \left\{
 |c_u|\sqrt{\frac{b_{q(u)}}{a_u}}:u\in S
 \right\}.
 \label{eq:tpc55-transversal-singular-values}
\end{equation}
\end{proposition}

\begin{proof}
Every vector in \eqref{eq:tpc55-kernel-vectors} is killed because its two
contributions are \(c_{u_0}c_u-c_uc_{u_0}=0\).  The vectors are independent
and their number is \(|F_y|-1\), which is the local nullity.  If a coordinate
restriction contains a zero coefficient or two directions from one active
fiber, it has a kernel.  Conversely, under \eqref{eq:tpc55-coordinate-injectivity}
the restricted matrix has one nonzero entry in each occupied row and column,
giving \eqref{eq:tpc55-transversal-singular-values}.
\end{proof}

This necessary-and-sufficient statement is why a collision-free coordinate
selection must retain at most one terminal atom from each complete output
fiber.  Merely bounding the fiber size does not make the unrestricted map
injective.

\subsection{Canonical quotient and the minimum-norm lift}

We use the standard Moore--Penrose inverse convention
\cite{BenIsraelGreville2003}.

\begin{theorem}[Exact pseudoinverse]
\label{thm:tpc55-pseudoinverse}
For every \(w\in\cH_b(\cY)\), the Moore--Penrose inverse is
\begin{equation}
 (T^\dagger w)_u
 =
 \begin{cases}
 \dfrac{\overline{c_u}}{a_uS_{q(u)}}w_{q(u)},
   &q(u)\in\cY_+,\\[6pt]
 0,&q(u)\notin\cY_+.
 \end{cases}
 \label{eq:tpc55-pseudoinverse-formula}
\end{equation}
Let \(Q_+\) denote the \(b\)-orthogonal projection onto the output
coordinates indexed by \(\cY_+\).  Then
\begin{align}
 TT^\dagger&=Q_+,
 &T^\dagger T&=P,
 \label{eq:tpc55-pseudoinverse-identities}
\end{align}
and, for every \(w\) supported on \(\cY_+\), \(T^\dagger w\) is the unique
minimum-\(a\)-norm solution of \(Tz=w\), with
\begin{equation}
 \inf_{Tz=w}\|z\|_a^2
 =\sum_{y\in\cY_+}\frac{|w_y|^2}{S_y}.
 \label{eq:tpc55-minimum-lift-energy}
\end{equation}
If \(\cY_+\ne\varnothing\), then
\begin{equation}
 \|T^\dagger\|_{b\to a}
 =\left(\min_{y\in\cY_+}\rho_y\right)^{-1}.
 \label{eq:tpc55-pseudoinverse-norm}
\end{equation}
If \(\cY_+=\varnothing\), then \(T=T^\dagger=0\) and
\(\|T^\dagger\|_{b\to a}=0\).
\end{theorem}

\begin{proof}
The formula is the inverse on each one-dimensional nonzero singular
direction and zero on both the domain kernel and the inactive output
coordinates.  Direct substitution gives
\eqref{eq:tpc55-pseudoinverse-identities}.  Its input energy on one output
coordinate is
\[
 \sum_{u\in F_y}a_u
 \left|\frac{\overline{c_u}}{a_uS_y}w_y\right|^2
 =\frac{|w_y|^2}{S_y}.
\]
Orthogonality of the fibers proves the energy and operator-norm formulas.
When there is no active fiber, both operators are zero, which gives the last
assertion.
\end{proof}

Thus the quotient by the collision kernel has a canonical realization.
Choosing one representative coordinate in each fiber is useful for support
selection, but it is not the minimum-norm quotient unless the fiber already
has one active coordinate.

\subsection{Restricted subspaces and the exact missing angle}

Assume throughout this subsection that \(\cY_+\ne\varnothing\).  If
\(\cY_+=\varnothing\), then \(T=0\), so no nonzero subspace can be a domain
of injectivity for \(T\).

Let
\begin{equation}
 \rho_{\min}=\min_{y\in\cY_+}\rho_y,
 \qquad
 \rho_{\max}=\max_{y\in\cY_+}\rho_y.
 \label{eq:tpc55-rho-extremes}
\end{equation}
For a nonzero subspace \(V\subseteq\cH_a(\cU)\), define
\begin{equation}
 \alpha(V)
 :=\inf_{\substack{z\in V\\\|z\|_a=1}}\|Pz\|_a,
 \qquad
 \gamma_T(V)
 :=\inf_{\substack{z\in V\\\|z\|_a=1}}\|Tz\|_b.
 \label{eq:tpc55-principal-angle}
\end{equation}

\begin{theorem}[Sharp angle criterion]
\label{thm:tpc55-angle-criterion}
For every nonzero subspace \(V\), in finite dimension,
\begin{equation}
 \alpha(V)>0
 \quad\Longleftrightarrow\quad
 V\cap\ker T=\{0\},
 \label{eq:tpc55-angle-injectivity}
\end{equation}
and
\begin{equation}
 \rho_{\min}\alpha(V)
 \le\gamma_T(V)
 \le\rho_{\max}\alpha(V).
 \label{eq:tpc55-angle-bounds}
\end{equation}
After normalizing the output coordinate \(y\) by \(1/\rho_y\), the
resulting operator \(\widehat T\) satisfies the exact identity
\begin{equation}
 \|\widehat Tz\|_b=\|Pz\|_a,
 \qquad
 \gamma_{\widehat T}(V)=\alpha(V).
 \label{eq:tpc55-normalized-angle-identity}
\end{equation}
\end{theorem}

\begin{proof}
By \cref{thm:tpc55-complete-svd},
\[
 \rho_{\min}^2\|Pz\|_a^2
 \le\|Tz\|_b^2
 \le\rho_{\max}^2\|Pz\|_a^2.
\]
Taking the infimum on the unit sphere of \(V\) proves
\eqref{eq:tpc55-angle-bounds}.  Compactness of that sphere gives
\eqref{eq:tpc55-angle-injectivity}.  Dividing each nonzero singular value
by itself gives \eqref{eq:tpc55-normalized-angle-identity}.
\end{proof}

Pointwise injectivity on a sequence of changing subspaces is therefore not
enough for a uniform theorem.  The subspaces must remain quantitatively
separated from the explicit collision kernel.

\subsection{Alignment of a distinguished vector}

For \(z\in\cH_a(\cU)\), put
\begin{equation}
 E_y(z):=\sum_{u\in F_y}a_u|z_u|^2,
 \qquad
 A_y(z):=\sum_{u\in F_y}c_uz_u,
 \label{eq:tpc55-fiber-energy-amplitude}
\end{equation}
and, when \(S_yE_y(z)>0\), define
\begin{equation}
 \eta_y(z):=\frac{|A_y(z)|^2}{S_yE_y(z)};
 \label{eq:tpc55-fiber-alignment}
\end{equation}
set \(\eta_y=0\) in the zero case.

\begin{theorem}[Exact visible-energy identity]
\label{thm:tpc55-visible-energy}
For every \(z\) and every active fiber \(y\in\cY_+\),
\begin{equation}
 0\le\eta_y(z)\le1,
 \qquad
 \|P_yz\|_a^2=\eta_y(z)E_y(z),
 \label{eq:tpc55-local-visible-energy}
\end{equation}
and
\begin{equation}
 \boxed{
 \|Tz\|_b^2
 =\sum_{y\in\cY_+}\rho_y^2\eta_y(z)E_y(z).}
 \label{eq:tpc55-exact-alignment-identity}
\end{equation}
On an active fiber, equality \(\eta_y=1\) holds precisely when the
restriction of \(z\) to that fiber is a nonzero scalar multiple of \(v_y\).
Equality \(\eta_y=0\) holds precisely when that restriction lies in the
fiber kernel, including the zero vector.  On an inactive fiber
\(y\notin\cY_+\), one has \(\eta_y=0\) by convention, all coefficients
\(c_u\) on \(F_y\) vanish, and the entire fiber is killed by \(T\).
\end{theorem}

\begin{proof}
Cauchy--Schwarz in the weighted fiber gives
\(|A_y|^2\le S_yE_y\).  Formula
\eqref{eq:tpc55-visible-projection} gives
\(\|P_yz\|_a^2=|A_y|^2/S_y\).  Multiplying by
\(\rho_y^2=b_yS_y\) and summing gives
\eqref{eq:tpc55-exact-alignment-identity}.  The equality cases are the
equality case of Cauchy--Schwarz and the definition of the kernel.
\end{proof}

For the physical terminal map, the inherited atomic norms have
\(a_u=b_y=1\) and
\begin{equation}
 c_u=\mu(r_u)\Lambda(p_ur_u).
 \label{eq:tpc55-physical-svd-coefficient}
\end{equation}
The theorem must be applied with these literal coefficients.  Choosing
artificial weights \(a_u=|c_u|^2\) would hide the terminal amplitude in a
new domain norm and is not a physical normalization theorem.
